\documentclass[11pt,a4paper]{report}

\usepackage[a4paper,margin=2.5cm]{geometry}
\usepackage[T1]{fontenc}
\usepackage[utf8]{inputenc}
\usepackage[italian]{babel}
\usepackage{lmodern}
\usepackage{microtype}
\usepackage{booktabs}
\usepackage{longtable}
\usepackage{tabularx}
\usepackage{array}
\usepackage{amsmath}
\usepackage{hyperref}
\usepackage{enumitem}
\usepackage{xcolor}
\usepackage{listings}
\usepackage{fancyhdr}

\hypersetup{
  colorlinks=true,
  linkcolor=blue!50!black,
  urlcolor=blue!50!black,
  pdftitle={Documentazione Tecnica regular\_clone},
  pdfauthor={Team Progetto}
}

\lstdefinelanguage{json}{
  basicstyle=\ttfamily\small,
  numbers=left,
  numberstyle=\scriptsize\color{gray},
  stepnumber=1,
  numbersep=8pt,
  showstringspaces=false,
  breaklines=true,
  frame=single,
  rulecolor=\color{black!20},
  backgroundcolor=\color{black!2},
}

\lstset{
  basicstyle=\ttfamily\small,
  frame=single,
  breaklines=true,
  backgroundcolor=\color{black!2},
  rulecolor=\color{black!20}
}

\pagestyle{fancy}
\fancyhf{}
\fancyhead[L]{regular\_clone}
\fancyhead[R]{Documentazione Tecnica}
\fancyfoot[C]{\thepage}
\setlength{\headheight}{14pt}

\setlist[itemize]{topsep=2pt,itemsep=2pt,parsep=0pt}
\setlist[enumerate]{topsep=2pt,itemsep=2pt,parsep=0pt}

\title{
  \vspace{2cm}
  \textbf{Documentazione Tecnica del Progetto regular\_clone}\\
  \large Allineamento codice-runtime e metodologia produttiva
}
\author{Team Progetto}
\date{23 Marzo 2026}

\begin{document}
\maketitle
\tableofcontents
\newpage

\chapter{Document Control}

\section{Revision History}
\begin{longtable}{@{}p{2.2cm}p{2.2cm}p{3.6cm}p{6.8cm}@{}}
\toprule
Versione & Data & Autore & Descrizione \\ \midrule
1.9 & 2026-03-23 & Team progetto & Introduzione baseline startup low-cost ufficiale con selezione min-CAPEX vincolata (\texttt{startup\_low\_capex + grid\_only\_retrofit + core\_efficiency\_extended}), quality floor nativo \texttt{62} su modalita resa, telemetria energetica opzionale \texttt{power\_*} nei control tick e aggiornamento runbook/documentazione business. \\
1.8 & 2026-03-23 & Team progetto & Aggiornata semantica resa annua da vincolo \texttt{<=} a target esatto \texttt{=} con pianificazione area attiva necessaria; aggiornati optimizer/realtime/planner/economics, test e memo business. \\
1.7 & 2026-03-23 & Team progetto & Ottimizzazione flusso editoriale: nuovo capitolo dedicato alle assunzioni sperimentali con tabella master (\texttt{hard runtime}/\texttt{calibrabile}), KPI economici conservativi su resa pianificata e inserimento memo business investitore in \texttt{docs/business}. \\
1.6 & 2026-03-23 & Team progetto & Inserita sezione esplicita sui valori sperimentali assunti pre-calibrazione (fenologia, coefficienti genetici/famiglia, bounds calibrazione twin) e riorganizzazione path repository (\texttt{docs/*}, \texttt{economics/*}). \\
1.5 & 2026-03-23 & Team progetto & Hardening pre-calibrazione: template dataset self-contained e family-coherent, rimozione variabili durata non piu ottimizzabili (\texttt{veg\_days}/\texttt{flower\_*\_days}), KPI economici estesi (ricavi, EBITDA, ROI, break-even) e script operativo \texttt{scripts/run\_economic\_analysis.py}. \\
1.4 & 2026-03-22 & Team progetto & Enforcement hard del cap energia in realtime anche con \texttt{--profile-json} (\texttt{controller\_literature\_realtime.py} e \texttt{core/realtime\_io.py}); aggiornamento sezione test/coerenza. \\
1.3 & 2026-03-22 & Team progetto & Introduzione vincolo hard \texttt{yield\_cap\_annual\_kg} in optimizer/realtime, planner default a 80 kg/anno, integrazione modulo economico locale \texttt{economics/cea\_economic\_analysis.py}. \\
1.2 & 2026-03-22 & Team progetto & Supporto nativo realtime per \texttt{max\_yield\_energy}/\texttt{max\_quality\_energy} e sezione esplicativa ad alto livello per stakeholder non tecnici. \\
1.1 & 2026-03-22 & Team progetto & Riorganizzazione capitolo fenologico: parametri \texttt{10--14 gg germinazione} e \texttt{27--35 gg seed->talea} dichiarati hard-coded con evidenza implementativa. \\
1.0 & 2026-03-22 & Team progetto & Prima emissione tecnica dedicata al runtime \texttt{regular\_clone}, con policy codice dominante e mappa scostamenti dal business plan. \\
\bottomrule
\end{longtable}

\section{Scopo del documento}
Questo documento definisce l'architettura tecnica, la formulazione tecnico-matematica e le procedure operative del progetto \texttt{regular\_clone}.\\
La fonte implementativa primaria e il codice del repository corrente (\texttt{regular\_clone/*}).\\
La fonte metodologica di riferimento e il documento interno \emph{BUISNESS PLAN (1)}.

\section{Perimetro e policy}
\begin{itemize}
  \item \textbf{Policy di allineamento}: codice dominante.
  \item \textbf{Uso del business plan}: riferimento metodologico per ciclo da talea, schema impiantistico DWC/NFT, gestione soluzione nutritiva e zoning operativo.
  \item \textbf{Scostamenti}: riportati esplicitamente nel Capitolo 14.
\end{itemize}

\newpage
\chapter{Executive Technical Overview}

\section{Obiettivo del sistema}
Il sistema implementa due anelli integrati:
\begin{enumerate}
  \item \textbf{Anello strategico offline}: ottimizzazione constrained Bayesian (\texttt{optimizer\_literature\_best.py}) su digital twin.
  \item \textbf{Anello operativo realtime}: controllo adattivo closed-loop (\texttt{controller\_literature\_realtime.py}) con quality gate, fallback di sicurezza e supervisione MPC opzionale.
\end{enumerate}

\section{Modalita operative}
Modalita supportate dall'optimizer e dal controller realtime:
\begin{itemize}
  \item \texttt{max\_yield}
  \item \texttt{max\_quality}
  \item \texttt{max\_yield\_energy}
  \item \texttt{max\_quality\_energy}
\end{itemize}
Le varianti \texttt{*\_energy} introducono un vincolo hard su energia per ciclo
(\texttt{energy\_cap\_kwh\_m2}, default 700.0, configurabile da CLI).\\
Tutte le modalita introducono inoltre un vincolo hard di resa annua farm-level
(\texttt{yield\_target\_annual\_kg}, alias CLI compatibile \texttt{yield\_cap\_annual\_kg}, default 80.0).\\
Per le modalita orientate alla resa (\texttt{max\_yield}, \texttt{max\_yield\_energy}) e attivo un guardrail
di qualita minimo (\texttt{quality\_floor}, default 62), verificato sia in ottimizzazione sia in realtime
anche in caso di profilo esterno caricato con \texttt{--profile-json}.\\
Dato il profilo ottimizzato, il runtime calcola l'area attiva pianificata necessaria per rispettare il target esatto:
\[
\texttt{planned\_active\_area\_m2\_for\_target} =
\frac{\texttt{yield\_target\_annual\_kg}}{\texttt{dry\_yield\_kg\_m2\_year}}
\]
e valida \(\texttt{planned\_active\_area\_m2\_for\_target} \le \texttt{farm\_active\_area\_m2}\).\\
Nel loop di controllo, le modalita \texttt{*\_energy} usano la stessa logica closed-loop della modalita base corrispondente
(\texttt{max\_yield} oppure \texttt{max\_quality}), ma operano su un profilo ottimizzato con vincolo energetico.
Se viene passato un profilo esterno (\texttt{--profile-json}), il runtime valida comunque in avvio che
\texttt{energy\_kwh\_m2\_cycle <= energy\_cap\_kwh\_m2} per le modalita \texttt{*\_energy}.

\section{Come funziona in modo semplice (per investitori)}
Il sistema puo essere letto come una piattaforma di gestione rischio-rendimento per una farm idroponica:
\begin{enumerate}
  \item si calibra il modello su una postazione pilota, per allineare simulazione e comportamento reale;
  \item si genera una ricetta ottimale di ciclo (resa o qualita), con eventuale tetto energetico compatibile col budget;
  \item si applica un target hard di produzione annua in uguaglianza per mantenere la farm entro il budget/obiettivo commerciale;
  \item in produzione i sensori misurano ogni giorno (o tick) lo stato reale della coltura e dell'ambiente;
  \item il controller confronta reale e target, calcola correzioni operative e invia comandi agli attuatori;
  \item ogni decisione e tracciata (DB + log), cosi da auditare performance, sicurezza e costi energetici.
\end{enumerate}
In termini economici, il valore e trasformare un processo agricolo complesso in una procedura controllabile,
misurabile e ripetibile, con compromesso esplicito tra produttivita e consumo energia.
Per una lettura business non tecnica, fare riferimento al memo investitore \texttt{docs/business/BUSINESS\_MEMO\_REGULAR\_CLONE\_IT.pdf}.

\section{KPI ufficiali}
\begin{longtable}{@{}p{4.0cm}p{4.0cm}p{5.8cm}@{}}
\toprule
KPI & Campo runtime & Significato \\ \midrule
Resa areale & \texttt{dry\_yield\_g\_m2} & Produzione secca per m2 a fine ciclo. \\
Qualita sintetica & \texttt{quality\_index} & Indice 0--100 derivato da segnali spettro/stress/salute. \\
Energia ciclo & \texttt{energy\_kwh\_m2} & Energia per m2 per ciclo (LED + HVAC proxy). \\
Efficienza energetica & \texttt{g\_per\_kwh} & Resa secca per kWh. \\
Rischio processo & \texttt{penalty} & Penalita da violazioni VPD, pH, EC, DO, RH. \\
Rischio fitosanitario & \texttt{disease\_pressure}, \texttt{hlvd\_pressure} & Proxy cumulativi sanitari. \\
Robustezza BO & \texttt{feasibility\_violation}, \texttt{uncertainty} & Coerenza vincoli e dispersione multi-run. \\
\bottomrule
\end{longtable}

\section{Allineamento metodologico di alto livello}
Elementi del business plan direttamente riflessi nel runtime:
\begin{itemize}
  \item ciclo produttivo da talea;
  \item \texttt{propagation\_system = DWC};
  \item \texttt{production\_system = NFT\_recirculating};
  \item fotoperiodicita controllata in fase di fioritura.
\end{itemize}

\newpage
\chapter{Assunzioni Sperimentali (Baseline Pre-Calibrazione)}
\label{chap:assunzioni_sperimentali}

\section{Policy assunzioni}
Prima della calibrazione su dati pilota, il runtime usa un set di assunzioni sperimentali codificate.
Ogni assunzione e classificata come:
\begin{itemize}
  \item \textbf{hard runtime}: vincolo non ottimizzabile nel loop giornaliero;
  \item \textbf{calibrabile}: parametro modificabile in fase di fit (\texttt{calibrate\_twin.py}).
\end{itemize}

\section{Tabella master assunzioni indica/sativa}
\begin{longtable}{@{}p{3.5cm}p{2.0cm}p{2.0cm}p{2.2cm}p{3.8cm}@{}}
\toprule
Parametro & Indica & Sativa & Stato & Fonte implementativa \\ \midrule
\texttt{clone\_cycle\_days\_from\_cutting} & 105 & 120 & hard runtime & \texttt{core/model.py} \\
\texttt{stage\_days\_from\_cutting} & 12/24/7/32/30 & 12/27/7/37/37 & hard runtime & \texttt{core/model.py} \\
\texttt{flower\_density\_pl\_m2} & 4.0 & 2.0 & hard runtime & \texttt{core/model.py}, \texttt{core/production\_planner.py} \\
\texttt{seed\_germination\_days} & 10--14 & 10--14 & hard runtime & \texttt{SEED\_GERMINATION\_DAYS\_RANGE} \\
\texttt{seed\_to\_cutting\_days} & 27--35 & 27--35 & hard runtime & metadata profilo regular \\
\texttt{yield\_potential\_scale} & 0.99 (fam.) & 1.02 (fam.) & calibrabile (indiretto) & \texttt{configs/genetic\_profiles.json} \\
\texttt{quality\_potential\_scale} & 0.98 (fam.) & 1.03 (fam.) & calibrabile (indiretto) & \texttt{configs/genetic\_profiles.json} \\
\texttt{disease\_pressure\_scale} & 0.97 (fam.) & 1.03 (fam.) & calibrabile (indiretto) & \texttt{configs/genetic\_profiles.json} \\
\texttt{hlvd\_pressure\_scale} & 1.00 (fam.) & 1.01 (fam.) & calibrabile (indiretto) & \texttt{configs/genetic\_profiles.json} \\
\texttt{TwinCalibration.yield\_gain\_scale} & 0.60--1.60 & 0.60--1.60 & calibrabile (diretto) & \texttt{calibrate\_twin.py} \\
\texttt{TwinCalibration.quality\_gain\_scale} & 0.60--1.60 & 0.60--1.60 & calibrabile (diretto) & \texttt{calibrate\_twin.py} \\
\texttt{TwinCalibration.energy\_scale} & 0.60--1.80 & 0.60--1.80 & calibrabile (diretto) & \texttt{calibrate\_twin.py} \\
\texttt{TwinCalibration.penalty\_scale} & 0.50--2.50 & 0.50--2.50 & calibrabile (diretto) & \texttt{calibrate\_twin.py} \\
\bottomrule
\end{longtable}

\section{Nota di lettura}
Questo capitolo e la baseline ufficiale per calibrazione e audit.
Il dettaglio matematico completo resta nei capitoli successivi (Parametrizzazione, Formulazione, Twin, BO).

\newpage
\chapter{Architettura Software}

\section{Layer applicativi}
\begin{longtable}{@{}p{2.8cm}p{5.0cm}p{7.2cm}@{}}
\toprule
Layer & File principali & Responsabilita \\ \midrule
Knowledge & \texttt{core/genetics.py}, \texttt{core/literature.py}, \texttt{configs/*.json} & Profili genetici, priors cultivar, registry fonti, override bounds/coefficienti. \\
Model & \texttt{core/model.py} & Dataclass dominio, builder profilo strategico, controller adattivo, digital twin, serializzazione profilo/outcome. \\
Optimization & \texttt{optimizer\_literature\_best.py} & Constrained BO multi-GP con restart robusti e stima incertezza ensemble. \\
Realtime & \texttt{controller\_literature\_realtime.py}, \texttt{core/realtime\_io.py}, \texttt{core/mpc\_supervisor.py} & Ingestione payload multi-source, quality gate, fallback, emissione setpoint e MPC opzionale. \\
Observability & \texttt{core/storage.py}, \texttt{core/alerting.py}, \texttt{dashboard\_timeseries.py} & Persistenza SQLite, alerting rolling-window, dashboard HTML. \\
Calibration & \texttt{calibrate\_twin.py} & Fit parametri del twin su cicli osservati (train/val, bootstrap, governance). \\
Planning & \texttt{core/production\_planner.py}, \texttt{scripts/plan\_clone\_logistics.py} & Pianificazione logistica multi-ciclo talee con vincoli hard. \\
Validation & \texttt{tests/test\_*.py}, \texttt{scripts/run\_tests.ps1} & Verifica funzionale unit/e2e, sicurezza realtime, vincoli regular, planner. \\
\bottomrule
\end{longtable}

\section{Data flow end-to-end}
\begin{enumerate}
  \item caricamento profilo genetico/cultivar e bounds;
  \item calibrazione twin opzionale (\texttt{calibrate\_twin.py});
  \item ottimizzazione profilo con BO constrained;
  \item persistenza run optimizer su \texttt{optimizer\_runs};
  \item ingestione realtime payload sensore;
  \item quality gate sensori e gestione fault;
  \item correzione euristica setpoint baseline;
  \item raffinamento MPC opzionale;
  \item emissione \texttt{control\_tick} su stdout/JSONL/sink attuatore;
  \item logging DB e dashboard.
\end{enumerate}

\section{Mappa zone capannone vs dominio software}
\begin{longtable}{@{}p{4.2cm}p{4.0cm}p{6.8cm}@{}}
\toprule
Zona business plan & Copertura nel runtime & Note \\ \midrule
Zona talee (DWC) & \texttt{stage = propagation} & Coerente: propagazione in ambiente separato e conservativo. \\
Zona vegetativa (NFT) & \texttt{stage = vegetative} & Coerente: setpoint vegetativi e gestione nutrizione/traspirazione. \\
Zona fioritura (NFT/contesto serra) & \texttt{stage = flower\_early/flower\_late} & Coerente sul dominio colturale; il dettaglio serra naturale non e simulato esplicitamente. \\
Laboratorio + dry room + magazzino & Non modellati in logica di controllo & Trattati come perimetro operativo extra-software. \\
\bottomrule
\end{longtable}

\newpage
\chapter{Parametrizzazione Fenologica e Spazio Decisionale (Dettaglio Operativo)}

\section{Fasi fenologiche}
Ordine runtime:
\[
\texttt{propagation} \rightarrow \texttt{vegetative} \rightarrow \texttt{transition} \rightarrow \texttt{flower\_early} \rightarrow \texttt{flower\_late}
\]

\section{Vincoli varietali e produttivi runtime}
Il profilo di default e \texttt{regular\_photoperiodic} con famiglia default \texttt{indica\_dominant}.\\
Il \texttt{StrategyBuilder} impone durate hard per famiglia.

\begin{longtable}{@{}p{3.2cm}p{1.5cm}p{1.5cm}p{1.5cm}p{1.7cm}p{1.5cm}p{2.2cm}@{}}
\toprule
Famiglia & Prop. & Veg. & Trans. & Flower early & Flower late & Densita fioritura \\ \midrule
\texttt{indica\_dominant} & 12 & 24 & 7 & 32 & 30 & 4.0 pl/m2 \\
\texttt{sativa\_dominant} & 12 & 27 & 7 & 37 & 37 & 2.0 pl/m2 \\
\texttt{hybrid} & 12 & 26 & 7 & 34 & 33 & 3.0 pl/m2 \\
\bottomrule
\end{longtable}

Nel metadata profilo sono registrati anche:
\begin{itemize}
  \item \texttt{business\_plan\_cuttings\_per\_plant = 5.0};
  \item \texttt{business\_plan\_no\_dedicated\_mother\_plants = True};
  \item finestra \texttt{seed\_to\_cutting} hard-coded: 27.0--35.0 giorni.
\end{itemize}

\section{Vincoli hard pre-ciclo (seme -> talea)}
I parametri pre-ciclo sono trattati come \textbf{costanti hard-coded} e non vengono ottimizzati dalla BO.
Essi sono usati per governance del profilo e tracciabilita metodologica.

\begin{longtable}{@{}p{4.2cm}p{2.0cm}p{3.2cm}p{4.0cm}@{}}
\toprule
Parametro & Valore hard & Implementazione & Impatto runtime \\ \midrule
\texttt{SEED\_GERMINATION\_DAYS\_RANGE} & 10--14 gg & \texttt{core/model.py} & Metadata profilo \texttt{business\_plan\_seed\_germination\_days\_min/max}. \\
\texttt{SEEDLING\_TO\_CUTTING\_DAYS\_RANGE} & 17--21 gg & \texttt{core/model.py} & Composizione finestra \texttt{seed->talea}. \\
\texttt{business\_plan\_seed\_to\_cutting\_days} & 27--35 gg & \texttt{core/model.py} & Tracciabilita tecnica del lead time pre-ciclo. \\
\bottomrule
\end{longtable}

Nota di perimetro: il ciclo simulato dal twin e dal controller realtime parte dalla talea distaccata
(\texttt{cutting\_cycle\_start = detached\_cutting}); la finestra \texttt{seed->talea} e quindi un vincolo hard pre-ciclo,
non un intervallo ottimizzabile nel loop giornaliero di controllo.

\section{Spazio parametrico ottimizzato (literature baseline)}
Bounds del costruttore \texttt{CannabisYieldLiteratureBuilder}:
\begin{longtable}{@{}p{3.4cm}p{1.2cm}p{1.2cm}p{2.0cm}p{5.0cm}@{}}
\toprule
Parametro & Min & Max & Unita & Fonte primaria runtime \\ \midrule
\texttt{veg\_ppfd} & 360 & 850 & umol m$^{-2}$ s$^{-1}$ & \texttt{rodriguez\_morrison\_2021\_light} \\
\texttt{flower\_ppfd} & 900 & 1800 & umol m$^{-2}$ s$^{-1}$ & \texttt{rodriguez\_morrison\_2021\_light} \\
\texttt{flower\_photoperiod\_h} & 12.0 & 14.0 & h & \texttt{ahrens\_2023\_is\_twelve\_hours}, \texttt{ahrens\_2024\_13h\_photoperiod} \\
\texttt{co2\_flower\_ppm} & 450 & 1200 & ppm & letteratura fisiologica CO2 \\
\texttt{air\_temp\_day\_c} & 23.0 & 29.0 & degC & letteratura CEA cannabis indoor \\
\texttt{rh\_day\_pct} & 50 & 70 & \% & modello VPD e rischio patogeni \\
\texttt{ec\_veg} & 1.4 & 2.4 & mS/cm & \texttt{saloner\_2020\_n} \\
\texttt{ec\_flower} & 1.8 & 3.0 & mS/cm & \texttt{saloner\_2019\_k}, \texttt{powell\_2026\_mass\_balance} \\
\texttt{ph\_target} & 5.7 & 6.2 & pH & \texttt{powell\_2026\_mass\_balance} \\
\texttt{n\_mg\_l} & 140 & 240 & mg/L & \texttt{saloner\_2020\_n} \\
\texttt{p\_mg\_l} & 15 & 60 & mg/L & registry nutrizione runtime \\
\texttt{k\_mg\_l} & 60 & 240 & mg/L & \texttt{saloner\_2019\_k} \\
\texttt{blue\_flower\_frac} & 0.10 & 0.28 & frazione & \texttt{holweg\_2024\_red\_white\_spectra} \\
\texttt{far\_red\_flower\_frac} & 0.02 & 0.14 & frazione & \texttt{holweg\_2024\_red\_white\_spectra} \\
\texttt{uvb\_late\_frac} & 0.00 & 0.10 & frazione & \texttt{llewellyn\_2023\_uv\_minimal\_effect} \\
\bottomrule
\end{longtable}
Nota: le durate fenologiche non sono piu variabili decisionali della BO; il runtime usa solo durate hard family-specific.

\section{Valori sperimentali assunti: rimando}
La baseline pre-calibrazione e formalizzata nel Capitolo~\ref{chap:assunzioni_sperimentali}.
In questo capitolo si riportano solo gli elementi necessari al dominio decisionale dell'optimizer e del controller.

\section{Override genetici/famiglia/cultivar}
Precedenza effettiva dei bounds:
\[
\texttt{literature\_bounds} \rightarrow \texttt{genetic\_bound\_overrides} \rightarrow \texttt{family\_bound\_overrides} \rightarrow \texttt{cultivar\_bound\_overrides}
\]
I coefficienti modello sono composti in modo moltiplicativo:
\[
c_k = c^{base}_k \cdot c^{family}_k \cdot c^{cultivar}_k
\]
con clamp numerici applicati nel twin (\texttt{yield\_potential\_scale}, \texttt{quality\_potential\_scale}, \texttt{disease\_pressure\_scale}, \texttt{hlvd\_pressure\_scale}, \texttt{nutrient\_window\_scale}).

\section{Evidenze bibliografiche in registry}
\begin{longtable}{@{}p{3.7cm}p{0.9cm}p{2.3cm}p{7.0cm}@{}}
\toprule
\texttt{source\_id} & Anno & Compat. regular & Citazione sintetica \\ \midrule
\texttt{knight\_2010\_scrog\_hydroponic} & 2010 & no & Knight et al., Forensic Sci Int. \\
\texttt{saloner\_2019\_k} & 2019 & yes & Saloner et al., Front Plant Sci. \\
\texttt{saloner\_2020\_n} & 2020 & yes & Saloner et al., Front Plant Sci. \\
\texttt{rodriguez\_morrison\_2021\_light} & 2021 & yes & Rodriguez-Morrison et al., Front Plant Sci. \\
\texttt{dowling\_2022\_photoperiod\_meta} & 2022 & yes & Dowling et al., Front Plant Sci. \\
\texttt{moher\_2021\_photoperiodic\_response\_in\_vitro} & 2021 & yes & Moher et al., HortScience. \\
\texttt{ahrens\_2023\_is\_twelve\_hours} & 2023 & yes & Ahrens et al., Plants. \\
\texttt{ahrens\_2024\_13h\_photoperiod} & 2024 & yes & Ahrens et al., Plants. \\
\texttt{llewellyn\_2023\_uv\_minimal\_effect} & 2021 & yes & Rodriguez-Morrison et al., UV-B. \\
\texttt{holweg\_2024\_red\_white\_spectra} & 2024 & yes & Holweg et al., Front Plant Sci. \\
\texttt{powell\_2026\_mass\_balance} & 2026 & yes & Powell et al., Front Plant Sci. \\
\texttt{pitman\_2020\_pythium\_report} & 2021 & yes & Pitman et al., Plant Disease. \\
\texttt{punja\_2025\_hlvd\_management} & 2025 & yes & Punja, Plants. \\
\texttt{ioannidis\_2022\_ssr\_micropropagation} & 2022 & yes & Ioannidis et al., Plants. \\
\bottomrule
\end{longtable}

\newpage
\chapter{Formulazione Tecnico-Matematica del Modello}

\section{Variabili derivate}
\subsection*{VPD}
\begin{equation}
\mathrm{VPD} = e_s(T_{air}) - e_a(T_{air}, RH)
\end{equation}
\begin{equation}
e_s = 0.6108\exp\left(\frac{17.27T}{T+237.3}\right), \quad e_a = e_s\frac{RH}{100}
\end{equation}

\subsection*{DLI}
\begin{equation}
\mathrm{DLI} = 0.0036 \cdot \mathrm{PPFD} \cdot \mathrm{photoperiod\_h}
\end{equation}

\section{Dinamica giornaliera}
Il twin usa aggiornamenti laggati con rumore gaussiano:
\begin{equation}
x_t = \mathrm{clip}(x_{t-1} + \alpha(x^\star - x_{t-1}) + \epsilon_t)
\end{equation}
dove \(\epsilon_t \sim \mathcal{N}(0,\sigma^2)\), con \(\alpha\) e \(\sigma\) specifici per variabile (\texttt{t\_air\_c}, \texttt{rh\_pct}, \texttt{co2\_ppm}, \texttt{ppfd}, \texttt{t\_solution\_c}, \texttt{ec\_ms\_cm}, \texttt{ph}, \texttt{do\_mg\_l}).

\section{Funzioni di risposta}
\begin{equation}
f_G(x)=\exp\left(-\frac{(x-\mu)^2}{2\sigma^2}\right), \qquad
f_L(x)=\frac{1}{1+\exp\left(-\frac{x-c}{s}\right)}
\end{equation}
Le funzioni gaussiane/logistiche sono usate per temperatura, VPD, nutrizione, root-zone, CO2 e spettro.

\section{Pesi fenologici e gain}
\begin{longtable}{@{}p{3.2cm}p{2.0cm}p{2.0cm}p{6.8cm}@{}}
\toprule
Fase & Peso resa & Peso qualita & Uso nel modello \\ \midrule
\texttt{propagation} & 0.06 & 0.00 & rooting e stabilizzazione iniziale \\
\texttt{vegetative} & 0.32 & 0.06 & espansione canopy \\
\texttt{transition} & 0.78 & 0.20 & transizione fisiologica \\
\texttt{flower\_early} & 1.28 & 0.75 & bulk resa fioritura \\
\texttt{flower\_late} & 1.06 & 1.15 & maturazione qualita \\
\bottomrule
\end{longtable}

\begin{equation}
\Delta Y_t = \mathrm{DLI}_t \cdot 0.52 \cdot w_Y(stage) \cdot canopy \cdot \Phi_{growth} \cdot health
\end{equation}
\begin{equation}
\Delta Q_t = w_Q(stage)\cdot \Phi_{spectrum,stress,root,health}
\end{equation}
con scaling da coefficienti genetici/famiglia/cultivar e da \texttt{TwinCalibration}.

\section{Penalty e rischio fitosanitario}
Incremento \texttt{disease\_pressure}:
\begin{equation}
\Delta D_t = 0.020\max(RH-72,0)+0.090\max(0.82-VPD,0)+0.070\max(T_{sol}-22,0)+0.095\max(7-DO,0)
\end{equation}
Incremento \texttt{hlvd\_pressure}:
\begin{equation}
\Delta H_t = \left(0.0018(1-s)+0.0008\Delta D_t\right)\cdot scale_{hlvd}
\end{equation}
dove \(s\) e \texttt{sanitation\_level}.\\
La penalty aggrega eccessi/deficit su VPD, pH, EC, DO e RH in fioritura.

\newpage
\chapter{Controller Adattivo Closed Loop}

\section{Target VPD fase-specifici}
\begin{longtable}{@{}p{3.4cm}p{3.0cm}p{7.4cm}@{}}
\toprule
Fase & Target VPD (kPa) & Effetto operativo \\ \midrule
\texttt{propagation} & 0.40--0.80 & minimizzazione stress idrico iniziale \\
\texttt{vegetative} & 0.80--1.20 & sostegno crescita vegetativa \\
\texttt{transition} & 1.00--1.35 & passaggio fenologico controllato \\
\texttt{flower\_early} & 1.10--1.45 & robustezza set infiorescenze \\
\texttt{flower\_late} & 1.00--1.35 & compromesso resa/qualita/rischio \\
\bottomrule
\end{longtable}

\section{Regole di controllo implementate}
Azioni principali in \texttt{AdaptiveController.adjust()}:
\begin{enumerate}
  \item \texttt{raise\_vpd} / \texttt{lower\_vpd}
  \item \texttt{dli\_recovery}
  \item \texttt{raise\_ph} / \texttt{lower\_ph}
  \item \texttt{raise\_ec} / \texttt{lower\_ec}
  \item \texttt{uptake\_balance\_high\_transpiration} / \texttt{uptake\_balance\_low\_transpiration}
  \item \texttt{cool\_root\_zone}, \texttt{raise\_do}
  \item \texttt{quality\_spectral\_shift} o \texttt{yield\_spectral\_shift} in \texttt{flower\_late}
\end{enumerate}

\section{Interazione con quality-gate e fallback}
Pipeline operativa per ciascun tick:
\begin{enumerate}
  \item parsing payload (\texttt{sensor\_payload\_to\_state});
  \item valutazione qualita sensore (\texttt{evaluate\_sensor\_quality});
  \item fallback safe su payload invalidi, stream error/idle, hard fault sensore;
  \item enforce limiti hard/rate (\texttt{enforce\_setpoint\_limits});
  \item MPC opzionale (\texttt{--use-mpc-supervisor});
  \item emissione \texttt{control\_tick}.
\end{enumerate}

\newpage
\chapter{Digital Twin: Dettagli Implementativi}

\section{Ciclo di simulazione}
\texttt{CEADigitalTwin.simulate\_cycle()}:
\begin{enumerate}
  \item inizializza stato da \texttt{propagation};
  \item determina fase attiva per giorno;
  \item applica controllo adattivo;
  \item avanza stato con \texttt{\_advance\_one\_day};
  \item accumula yield/quality/energy/penalty;
  \item applica correttivi finali malattia/HLVd + \texttt{TwinCalibration} post-ciclo.
\end{enumerate}

\section{Robustezza numerica e guardrail}
\begin{itemize}
  \item clipping sistematico su variabili fisiche;
  \item seed espliciti per riproducibilita;
  \item clamp parametri calibrazione (\texttt{clamp\_twin\_calibration});
  \item normalizzazione spettro (\texttt{\_normalize\_spectrum}).
\end{itemize}

\section{Parametri di calibrazione twin}
\begin{longtable}{@{}p{4.0cm}p{2.2cm}p{2.2cm}p{5.0cm}@{}}
\toprule
Parametro & Min & Max & Ruolo \\ \midrule
\texttt{yield\_gain\_scale} & 0.60 & 1.60 & scaling gain giornaliero resa \\
\texttt{quality\_gain\_scale} & 0.60 & 1.60 & scaling gain giornaliero qualita \\
\texttt{transpiration\_scale} & 0.60 & 1.60 & scaling proxy traspirazione \\
\texttt{disease\_inc\_scale} & 0.50 & 2.00 & scala crescita pressione patogena \\
\texttt{hlvd\_inc\_scale} & 0.50 & 2.00 & scala crescita pressione HLVd \\
\texttt{penalty\_scale} & 0.50 & 2.50 & scala penalty di processo \\
\texttt{energy\_scale} & 0.60 & 1.80 & scala consumo energetico \\
\texttt{yield\_post\_scale} & 0.70 & 1.30 & scala resa post-ciclo \\
\texttt{quality\_post\_scale} & 0.70 & 1.30 & scala qualita post-ciclo \\
\texttt{quality\_post\_offset} & -20.0 & 20.0 & offset qualita post-ciclo \\
\bottomrule
\end{longtable}

\section{Output strutturato}
Output nativo:
\begin{itemize}
  \item \texttt{CycleOutcome} (KPI aggregati + log giornalieri);
  \item \texttt{outcome\_to\_dict()} con \texttt{adjustment\_summary} e \texttt{last\_days};
  \item \texttt{clone\_cycle\_derived\_metrics()} con KPI annualizzati.
\end{itemize}

\newpage
\chapter{Constrained Bayesian Optimization}

\section{Obiettivi e vincoli}
Per \texttt{max\_yield} l'obiettivo primario e \(\max\ \texttt{dry\_yield\_g\_m2}\).\\
Per \texttt{max\_quality} l'obiettivo primario e \(\max\ \texttt{quality\_index}\) con floor resa \texttt{quality\_y\_min}.\\
Le varianti \texttt{*\_energy} mantengono l'obiettivo base e aggiungono vincolo hard energia.
In tutte le modalita e attivo il target annuo farm-level \texttt{yield\_target\_annual\_kg}.
L'ottimizzazione richiede una resa areale annua minima \texttt{yield\_target\_kg\_m2\_year}
per rendere il target raggiungibile con l'area attiva disponibile.

\section{Violation function}
\texttt{\_literature\_violation()} penalizza:
\begin{itemize}
  \item fioritura sottodimensionata vs fase vegetativa;
  \item fotoperiodo fuori finestra 12--14h;
  \item eccessi di P/EC;
  \item criticita processo/sanitarie via \texttt{penalty}, \texttt{disease\_pressure}, \texttt{hlvd\_pressure}.
\end{itemize}

\section{Modello probabilistico multi-GP}
Modelli Gaussian Process allenati:
\begin{itemize}
  \item \texttt{gp\_obj}
  \item \texttt{gp\_yield}
  \item \texttt{gp\_violation}
  \item \texttt{gp\_penalty}
  \item \texttt{gp\_disease}
  \item \texttt{gp\_hlvd}
  \item \texttt{gp\_energy}
\end{itemize}
Kernel comune:
\[
K = \texttt{ConstantKernel} \times \texttt{Matern}(\nu=2.5) + \texttt{WhiteKernel}
\]

\section{Acquisition e fattibilita}
Acquisition:
\[
A(x) = EI(x)\cdot P_{feasible}(x)
\]
dove:
\[
P_{feasible}(x)=P(v\le 0.50)\cdot P(pen\le 15)\cdot P(dis\le 3)\cdot P(hlvd\le 0.08)\cdot P(energy\le cap)\cdot P(yield\ge floor)
\]
I fattori energia/floor resa sono attivati in base alla modalita.

\section{Strategia restart e robustezza}
\begin{itemize}
  \item \texttt{max\_yield}: restart multipli + robust score Monte Carlo multi-seed.
  \item \texttt{max\_quality}: restart multipli fino a candidato fattibile.
  \item incertezza finale: ensemble run con statistiche \texttt{p05/p50/p95}, \texttt{std}, \texttt{feasible\_rate}.
\end{itemize}

\newpage
\chapter{Contratto Dati Realtime e Integrazione I/O}

\section{Campi obbligatori e opzionali payload sensore}
\begin{longtable}{@{}p{3.6cm}p{2.8cm}p{6.8cm}@{}}
\toprule
Campo & Tipo & Stato \\ \midrule
\texttt{t\_air\_c} & float & obbligatorio \\
\texttt{rh\_pct} & float & obbligatorio \\
\texttt{co2\_ppm} & float & obbligatorio \\
\texttt{ppfd} & float & obbligatorio \\
\texttt{t\_solution\_c} & float & obbligatorio \\
\texttt{do\_mg\_l} & float & obbligatorio \\
\texttt{ec\_ms\_cm} & float & obbligatorio \\
\texttt{ph} & float & obbligatorio \\
\texttt{timestamp} & ISO string & opzionale \\
\texttt{transpiration\_l\_m2\_day} & float & opzionale \\
\texttt{dli\_prev} & float & opzionale \\
\texttt{vpd\_kpa} & float & opzionale \\
\texttt{disease\_pressure} & float & opzionale \\
\texttt{hlvd\_pressure} & float & opzionale \\
\texttt{power\_total\_kw} & float & opzionale (telemetria energetica) \\
\texttt{power\_led\_kw} & float & opzionale (telemetria energetica) \\
\texttt{power\_hvac\_kw} & float & opzionale (telemetria energetica) \\
\texttt{power\_pumps\_kw} & float & opzionale (telemetria energetica) \\
\bottomrule
\end{longtable}

\section{Sorgenti e sink supportati}
Sorgenti:
\begin{itemize}
  \item \texttt{jsonl\_file}
  \item \texttt{stdin\_json}
  \item \texttt{http\_poll}
  \item \texttt{serial\_json}
  \item \texttt{mock\_stream}
\end{itemize}
Sink:
\begin{itemize}
  \item stdout JSON;
  \item append su JSONL;
  \item HTTP POST attuatore;
  \item seriale attuatore.
\end{itemize}

\section{Schema output \texttt{control\_tick}}
\begin{longtable}{@{}p{3.8cm}p{9.4cm}@{}}
\toprule
Campo & Descrizione \\ \midrule
\texttt{event} & Evento strutturato (\texttt{control\_tick}). \\
\texttt{ts} & Timestamp ISO. \\
\texttt{mode}, \texttt{cycle\_day}, \texttt{stage} & Contesto runtime. \\
\texttt{actions} & Lista azioni controller/fallback/MPC. \\
\texttt{sensor} & Stato sensore validato. \\
\texttt{baseline\_setpoint} & Setpoint nominale stage. \\
\texttt{heuristic\_setpoint} & Output controller euristico pre-limiti/MPC. \\
\texttt{recommended\_setpoint} & Setpoint finale emesso. \\
\texttt{sensor\_quality} & Score + warning + hard fault. \\
\texttt{energy\_observability\_kw} & Telemetria potenza opzionale estratta da \texttt{power\_*} del payload. \\
\texttt{mpc} & Diagnostica MPC (se attivo). \\
\texttt{fault} & Oggetto fault in fallback ticks. \\
\texttt{governance} & \texttt{run\_id}, \texttt{profile\_signature}. \\
\bottomrule
\end{longtable}

\newpage
\chapter{Runbook Operativo}

\section{Setup ambiente}
\begin{lstlisting}
cd regular_clone
python -m venv .venv
.\.venv\Scripts\python.exe -m pip install --upgrade pip
.\.venv\Scripts\python.exe -m pip install -r requirements.txt
.\.venv\Scripts\python.exe -m pip install -r requirements_realtime_optional.txt
\end{lstlisting}

\section{Calibrazione twin}
\begin{lstlisting}
.\.venv\Scripts\python.exe .\calibrate_twin.py ^
  --dataset-json .\calibration_dataset_template.json ^
  --out-json .\runtime\profiles\twin_calibration.json ^
  --store-db .\runtime\db\cea_timeseries.db
\end{lstlisting}
Il template e self-contained (profili inline) e validato con controlli hard su coerenza famiglia/densita/sottosistema idroponico.

\section{Ottimizzazione offline}
\begin{lstlisting}
.\.venv\Scripts\python.exe .\optimizer_literature_best.py ^
  --mode both_energy --quality-floor 62 --energy-cap-kwh-m2 700 --json-only
\end{lstlisting}

\section{Realtime}
\textbf{Con auto-generazione profilo}:
\begin{lstlisting}
.\.venv\Scripts\python.exe .\controller_literature_realtime.py ^
  --mode max_yield --quality-floor 62 --source serial_json --serial-port COM4 --serial-baud 115200 ^
  --watchdog-timeout-s 5 --yield-target-annual-kg 80 --farm-active-area-m2 120 ^
  --store-db .\runtime\db\cea_timeseries.db
\end{lstlisting}

\textbf{Con vincolo energia nativo in realtime}:
\begin{lstlisting}
.\.venv\Scripts\python.exe .\controller_literature_realtime.py ^
  --mode max_yield_energy --quality-floor 62 --energy-cap-kwh-m2 640 ^
  --yield-target-annual-kg 80 --farm-active-area-m2 120 ^
  --source serial_json --serial-port COM4 --serial-baud 115200 ^
  --start-date 2026-04-01
\end{lstlisting}

\textbf{Con profilo gia esportato + MPC}:
\begin{lstlisting}
.\.venv\Scripts\python.exe .\controller_literature_realtime.py ^
  --mode max_quality --profile-json .\runtime\profiles\literature_best_profile_runtime.json ^
  --source mock_stream --use-mpc-supervisor --mpc-horizon 6 --mpc-candidates 96
\end{lstlisting}

\section{Analisi economica pre-go-live}
\begin{lstlisting}
.\.venv\Scripts\python.exe .\scripts\run_economic_analysis.py ^
  --target-annual-kg 80 --price-eur-g 4.0 ^
  --mix-indica 0.5 --mix-sativa 0.5 ^
  --yield-source optimizer --optimizer-mode max_yield_energy ^
  --quality-floor 62 ^
  --infrastructure-profile startup_low_capex ^
  --energy-architecture grid_only_retrofit ^
  --monitoring-tier core_efficiency_extended ^
  --energy-cap-kwh-m2-cycle 700 ^
  --economic-yield-basis-policy target ^
  --out-json .\runtime\logs\economic_analysis_project.json
\end{lstlisting}
Output principale: KPI annui (ricavi, EBITDA, ROI, payback, break-even) con base economica esplicita.
Campi chiave: \texttt{target\_annual\_kg}, \texttt{plan\_projected\_annual\_yield\_kg},
\texttt{economic\_yield\_basis\_policy}, \texttt{economic\_yield\_basis\_kg},
\texttt{target\_vs\_planned\_kg}, \texttt{selected\_configuration},
\texttt{candidate\_rankings}, \texttt{constraints\_satisfied}.

\section{Baseline startup low-cost ufficiale}
Per lo stato attuale pre-calibrazione la baseline economico-tecnica ufficiale e:
\texttt{startup\_low\_capex + grid\_only\_retrofit + core\_efficiency\_extended}.\\
Il selettore economico valuta il catalogo completo e seleziona la prima configurazione
con \texttt{constraints\_satisfied.all = true} e \texttt{capex\_total} minimo.

\begin{longtable}{@{}p{5.0cm}p{7.0cm}@{}}
\toprule
Elemento & Politica baseline \\ \midrule
Energia & rete elettrica + predisposizione retrofit (PV/batteria in fase 2) \\
Monitoraggio & 8 sensori core obbligatori + telemetria energetica opzionale \texttt{power\_*} \\
Controllo resa & target annuo in uguaglianza (\texttt{yield\_target\_annual\_kg = 80} default) \\
Controllo energia & cap hard su modalita \texttt{*\_energy} \\
Controllo qualita & \texttt{quality\_floor = 62} su modalita resa \\
Vincoli fenologici & invariati (durate/densita hard indica/sativa, DWC/NFT) \\
\bottomrule
\end{longtable}

\section{Artefatti runtime}
\begin{longtable}{@{}p{4.2cm}p{7.8cm}@{}}
\toprule
Artefatto & Path default \\ \midrule
DB SQLite & \texttt{runtime/db/cea\_timeseries.db} \\
Output JSONL controllo & \texttt{runtime/jsonl/control\_output\_literature\_realtime.jsonl} \\
Profili runtime & \texttt{runtime/profiles/} \\
Dashboard HTML & \texttt{runtime/logs/dashboard\_timeseries.html} \\
Analisi economica JSON & \texttt{runtime/logs/economic\_analysis\_project.json} \\
\bottomrule
\end{longtable}

\newpage
\chapter{Validazione e Test}

\section{Strategia test}
La validazione copre:
\begin{itemize}
  \item unit test su modelli, bounds, qualita sensori, alerting;
  \item test e2e optimizer e calibrazione;
  \item test safety realtime (fallback + limiti hard/rate);
  \item test planner logistico multi-ciclo.
\end{itemize}

\section{Suite test corrente}
Comando ufficiale:
\begin{lstlisting}
.\.venv\Scripts\python.exe -m unittest discover -s tests -p "test_*.py" -v
\end{lstlisting}

\begin{longtable}{@{}p{5.2cm}p{8.0cm}@{}}
\toprule
File test & Copertura \\ \midrule
\texttt{test\_e2e\_optimizer.py} & Run modes optimizer, metriche derivate clone cycle, varianti energy, guardrail \texttt{quality\_floor}. \\
\texttt{test\_e2e\_calibration.py} & E2E CLI calibrazione twin su dataset sintetico. \\
\texttt{test\_realtime\_safety.py} & Fallback setpoint, emissione fallback tick, limiti hard/rate, parsing telemetria \texttt{power\_*}. \\
\texttt{test\_realtime\_stream\_events.py} & Eventi stream \texttt{idle/error/valid}. \\
\texttt{test\_sensor\_quality.py} & Range hard, jump hard/warning, scoring. \\
\texttt{test\_alerting.py} & Trigger/cooldown alert rolling-window. \\
\texttt{test\_model\_coefficients.py} & Composizione coefficienti base/famiglia/cultivar. \\
\texttt{test\_literature\_bounds\_and\_sources.py} & Integrita bounds e registry fonti. \\
\texttt{test\_reference\_audit\_summary.py} & Sintesi compatibilita referenze per profilo. \\
\texttt{test\_realtime\_mode\_support.py} & Coerenza mapping modalita energy in realtime + validazione parsing cap energia e helper runtime. \\
\texttt{test\_regular\_profile\_constraints.py} & Verifica durate/densita hard regular per famiglia e finestra hard \texttt{germinazione/seed->talea}. \\
\texttt{test\_production\_planner.py} & Vincoli planner 80 kg/anno, check saturazione e match esatto del target annuo da schedule settimanale. \\
\texttt{test\_economic\_hybrid\_model.py} & Verifica modulo economico locale (catalogo configurazioni, selettore min-CAPEX e backward compatibility legacy). \\
\bottomrule
\end{longtable}

\section{Criteri minimi di accettazione}
\begin{enumerate}
  \item test suite passata al 100\%;
  \item nessun hard fault sensore non gestito;
  \item fallback e alert persistiti in DB quando \texttt{--store-db} e attivo;
  \item coerenza KPI optimizer/realtime con vincoli modalita.
\end{enumerate}

\newpage
\chapter{Sicurezza Operativa}

\section{Guardrail minimi}
\begin{longtable}{@{}p{4.6cm}p{8.6cm}@{}}
\toprule
Guardrail & Implementazione \\ \midrule
Fallback recipe conservativo & \texttt{safe\_fallback\_setpoint()} e \texttt{emit\_safe\_fallback\_tick()}. \\
Hard limits setpoint & \texttt{enforce\_setpoint\_limits()} con bounds fisici per clima/nutrizione/spettro. \\
Rate limits per tick & Clamp variazione per \texttt{ppfd}, \texttt{t\_air\_c}, \texttt{rh\_pct}, \texttt{co2\_ppm}, \texttt{t\_solution\_c}, \texttt{do\_mg\_l}, \texttt{ec\_ms\_cm}, \texttt{ph}. \\
Quality gate sensori & \texttt{evaluate\_sensor\_quality()} con hard faults e warnings. \\
Watchdog stream & Trigger fallback su timeout (\texttt{--watchdog-timeout-s}). \\
\bottomrule
\end{longtable}

\section{Fault handling e policy fallback}
Eventi fault principali:
\begin{itemize}
  \item \texttt{safe\_fallback\_watchdog\_timeout}
  \item \texttt{safe\_fallback\_source\_error}
  \item \texttt{safe\_fallback\_sensor\_payload\_error}
  \item \texttt{safe\_fallback\_sensor\_quality\_fault}
\end{itemize}
Ogni fallback genera un \texttt{control\_tick} strutturato con campo \texttt{fault}.

\section{Alerting rolling-window}
\begin{longtable}{@{}p{4.0cm}p{2.3cm}p{2.3cm}p{4.6cm}@{}}
\toprule
Parametro & Default & Campo CLI & Significato \\ \midrule
Finestra tick & 60 & \texttt{--alert-window-ticks} & ampiezza finestra mobile \\
Soglia fallback rate & 0.15 & \texttt{--alert-fallback-rate-threshold} & trigger \texttt{fallback\_rate\_high} \\
Soglia timeout count & 3 & \texttt{--alert-timeout-count-threshold} & trigger \texttt{timeout\_burst} \\
Soglia payload errors & 3 & \texttt{--alert-payload-error-count-threshold} & trigger \texttt{payload\_error\_burst} \\
Soglia quality faults & 3 & \texttt{--alert-quality-fault-count-threshold} & trigger \texttt{sensor\_quality\_fault\_burst} \\
Cooldown & 30 & \texttt{--alert-cooldown-ticks} & anti-ripetizione alert \\
\bottomrule
\end{longtable}

\newpage
\chapter{Traceability e Governance}

\section{Matrice codice -> sezioni}
\begin{longtable}{@{}p{3.4cm}p{5.8cm}p{5.2cm}@{}}
\toprule
Dominio & File fonte di verita & Sezioni documento \\ \midrule
Architettura complessiva & \texttt{README.md}, entrypoint root, \texttt{core/*.py} & Cap. 2--3 \\
Fenologia e bounds & \texttt{core/model.py}, \texttt{optimizer\_literature\_best.py}, \texttt{configs/genetic\_profiles.json} & Cap. 4--5 \\
Controller realtime & \texttt{controller\_literature\_realtime.py}, \texttt{core/realtime\_io.py} & Cap. 6, 9, 12 \\
MPC supervisor & \texttt{core/mpc\_supervisor.py} & Cap. 6, 12 \\
BO vincolata & \texttt{optimizer\_literature\_best.py} & Cap. 8 \\
Modalita realtime energy-constrained & \texttt{controller\_literature\_realtime.py}, \texttt{core/realtime\_io.py}, \texttt{core/mpc\_supervisor.py} & Cap. 2, 6, 10 \\
Calibrazione twin & \texttt{calibrate\_twin.py} & Cap. 7, 10 \\
Storage e audit & \texttt{core/storage.py}, \texttt{core/governance.py} & Cap. 10, 13 \\
Planner logistico & \texttt{core/production\_planner.py} & Cap. 10, 14 \\
Blocco economico & \texttt{economics/cea\_economic\_analysis.py} & Cap. 2, 11, 15 \\
Validazione & \texttt{tests/test\_*.py}, \texttt{scripts/run\_tests.ps1} & Cap. 11 \\
\bottomrule
\end{longtable}

\section{Run identity e firme}
La governance runtime include:
\begin{itemize}
  \item \texttt{run\_id} via \texttt{new\_run\_id(prefix)};
  \item \texttt{git\_commit} via \texttt{git\_commit\_short()};
  \item firma payload via \texttt{dict\_signature()} (profili/run/runtime);
  \item firma file via \texttt{file\_signature()} usata nel flusso di calibrazione dataset.
\end{itemize}

\section{Persistenza auditabile}
Schema SQLite inizializzato da \texttt{init\_storage()}:
\begin{itemize}
  \item \texttt{control\_ticks}
  \item \texttt{optimizer\_runs}
  \item \texttt{twin\_calibrations}
  \item \texttt{alert\_events}
\end{itemize}
Ogni tabella mantiene \texttt{payload\_json} completo per audit retrospettivo.

\newpage
\chapter{Allineamento Metodologico al Business Plan}

\section{Trasposizione metodologia produttiva}
Mappa sintetica business plan \(\rightarrow\) runtime:
\begin{longtable}{@{}p{5.0cm}p{3.4cm}p{4.8cm}@{}}
\toprule
Elemento business plan & Stato runtime & Evidenza \\ \midrule
Ciclo da talea & implementato & stage \texttt{propagation} + metadata \texttt{cuttings\_source}. \\
DWC in radicazione & implementato & \texttt{propagation\_system = DWC}. \\
NFT in vegetativa/fioritura con ricircolo & implementato & \texttt{production\_system = NFT\_recirculating}. \\
Gestione nutrizione (EC/pH/NPK) & implementato & setpoint stage + regole controller + penalty model. \\
Zonizzazione capannone (talee/veg/fioritura) & implementato a livello logico & mapping stage fenologici. \\
Dry room/laboratorio/magazzino & non modellato nel controllo & perimetro operativo extra-software. \\
\bottomrule
\end{longtable}

\section{Scostamenti espliciti BP vs codice}
\begin{longtable}{@{}p{3.4cm}p{5.0cm}p{5.0cm}@{}}
\toprule
Tema & Business plan & Runtime \texttt{regular\_clone} \\ \midrule
Piante madri dedicate & previste in avvio processo & \texttt{business\_plan\_no\_dedicated\_mother\_plants=True} e source talee da piante in vegetativa. \\
Densita impianto & riferimenti testuali 15/9 (e passaggi 16 e 6/9) pl/m2 & vincoli hard profilo regular: indica 4, sativa 2, hybrid 3 pl/m2 in fioritura. \\
Durate fasi iniziali & radicazione circa 15 gg + vegetativa circa 25 gg & vincoli hard pre-ciclo: germinazione 10--14 gg, seed->talea 27--35 gg; da talea a raccolta: 12/24/7/32/30 (indica), 12/27/7/37/37 (sativa). \\
Contesto serra luce naturale & descritto per risparmio energetico & non modellato come forcing meteorologico/stagionale esplicito. \\
Zonizzazione non coltiva & dry room/laboratorio/magazzino dettagliati & non presenti nel dominio software di controllo/ottimizzazione. \\
\bottomrule
\end{longtable}

\section{Gap metodologici dichiarati}
\textbf{GAP-01}: il business plan non fornisce setpoint machine-ready completi per ogni fase (T/RH/CO2/EC/pH/DO).\\
\textbf{GAP-02}: il business plan non formalizza equazioni digital twin/BO/MPC; la formulazione deriva dal codice implementativo.

\section{Implicazioni ingegneristiche degli scostamenti}
\begin{itemize}
  \item il runtime e progettato per robustezza software e controllo ripetibile, non per replica letterale della BOM impiantistica;
  \item le assunzioni hard su ciclo/densita influenzano direttamente KPI annualizzati e piano logistica talee;
  \item la mancata modellazione meteo-serra suggerisce una roadmap di estensione verso forcing esogeno.
\end{itemize}

\newpage
\chapter{Limiti Noti e Roadmap}

\section{Limiti attuali}
\begin{itemize}
  \item modello digital twin semplificato (proxy biologici) e non gemello fisico completo;
  \item MPC candidate-based campionato (non NMPC continuo con solver dedicato);
  \item assenza modellazione esplicita condizioni meteo/fotoperiodo locale esterno;
  \item dipendenza dalla qualita dati nel loop realtime e dataset calibrazione.
\end{itemize}

\section{Roadmap tecnica suggerita}
\begin{enumerate}
  \item integrazione forcing meteo/fotoperiodo esterno nel twin e nel controller;
  \item estensione multi-sito/multi-lotto per calibrazione;
  \item stima online di stato/parametri (state estimation);
  \item estensione economico-energetica da caso statico a co-ottimizzazione operativa (oggi il modulo economico e locale ma ancora separato dal loop di controllo).
\end{enumerate}

\newpage
\chapter{Riferimenti Tecnici}

\section{Fonti runtime e interne}
\begin{enumerate}
  \item Team Progetto CEA (2026). \emph{BUISNESS PLAN (1)}. Documento interno metodologico.
  \item \texttt{regular\_clone/configs/literature\_sources.json}. Registry fonti con campo \texttt{compatibility\_regular}.
  \item \texttt{regular\_clone/configs/genetic\_profiles.json}. Profili genetici e override bounds/coefficienti.
  \item \texttt{regular\_clone/configs/cultivar\_catalog.json}. Priors cultivar e riferimenti associati.
\end{enumerate}

\section{Letteratura agronomica in registry}
\begin{enumerate}
  \item Knight G. et al. (2010), Forensic Sci Int. 202(1-3):36-44.
  \item Saloner A. et al. (2019), Front Plant Sci. 10:1369.
  \item Saloner A. et al. (2020), Front Plant Sci. 11:572293.
  \item Rodriguez-Morrison V. et al. (2021), Front Plant Sci. 12:646020.
  \item Rodriguez-Morrison V. et al. (2021), Front Plant Sci. 12:725078.
  \item Dowling A.V. et al. (2022), Front Plant Sci. 12:797425.
  \item Moher M. et al. (2021), HortScience 56(2):183-189.
  \item Ahrens A. et al. (2023), Plants 12(14):2605.
  \item Ahrens C.W. et al. (2024), Plants 13(3):433.
  \item Holweg M.M.S.F. et al. (2024), Front Plant Sci. 15:1393803.
  \item Powell K. et al. (2026), Front Plant Sci. 16:1753553.
  \item Pitman T.L. et al. (2021), Plant Disease 105(4):1230.
  \item Punja Z.K. (2025), Plants 14(5):830.
  \item Ioannidis K. et al. (2022), Plants 11(19):2569.
\end{enumerate}

\section{Riferimenti metodologici software/matematici}
\begin{enumerate}
  \item Allen R.G. et al. (1998). FAO56 crop evapotranspiration.
  \item Jones D.R. et al. (1998). Efficient Global Optimization.
  \item Gelbart M.A. et al. (2014). Bayesian optimization with unknown constraints.
  \item Rasmussen C.E., Williams C.K.I. (2006). Gaussian Processes for Machine Learning.
  \item Mayne D.Q. (2014). Model Predictive Control.
\end{enumerate}

\appendix

\chapter{Esempio payload sensore}
\begin{lstlisting}[language=json]
{
  "timestamp": "2026-03-22T10:15:00",
  "t_air_c": 25.8,
  "rh_pct": 61.0,
  "co2_ppm": 780.0,
  "ppfd": 930.0,
  "t_solution_c": 20.1,
  "do_mg_l": 7.5,
  "ec_ms_cm": 2.0,
  "ph": 5.9,
  "transpiration_l_m2_day": 3.2,
  "dli_prev": 38.5
}
\end{lstlisting}

\chapter{Esempio output control\_tick}
\begin{lstlisting}[language=json]
{
  "event": "control_tick",
  "ts": "2026-03-22T10:15:01",
  "mode": "max_yield",
  "cycle_day": 44,
  "stage": "flower_early",
  "actions": [
    "dli_recovery",
    "uptake_balance_high_transpiration",
    "enforce_limits",
    "mpc_supervisor"
  ],
  "sensor_quality": {
    "score": 94.0,
    "warnings": [],
    "hard_faults": []
  },
  "governance": {
    "run_id": "rt_20260322101501_ab12cd34",
    "profile_signature": "sha256..."
  }
}
\end{lstlisting}

\chapter{Requisiti software}
\begin{itemize}
  \item Python 3.11+
  \item \texttt{numpy >= 2.0.0}
  \item \texttt{scipy >= 1.7.0}
  \item \texttt{scikit-learn >= 1.0.0}
  \item opzionale realtime seriale: \texttt{pyserial >= 3.5}
\end{itemize}

\begin{lstlisting}
.\.venv\Scripts\python.exe -m pip install -r requirements.txt
.\.venv\Scripts\python.exe -m pip install -r requirements_realtime_optional.txt
\end{lstlisting}

\end{document}
